\documentclass[11pt]{article}
\usepackage[margin=0.9in]{geometry}
\usepackage{amsmath,amsfonts,amssymb,amsthm}
\usepackage{enumitem}
\usepackage[]{graphicx}
\usepackage{color,subfigure}
\usepackage{multicol}
\usepackage{float}
\usepackage[all]{xypic}
\usepackage[colorlinks=true,citecolor=cyan,linkcolor=magenta]{hyperref}
\usepackage{colonequals}

\usepackage{stmaryrd}
\usepackage{fancyhdr, lastpage}
\pagestyle{fancy}
%\fancyfoot[C]{{\thepage} of \pageref{LastPage}}



\DeclareMathOperator{\mSpec}{mSpec}
\DeclareMathOperator{\Spec}{Spec}
\DeclareMathOperator{\Ass}{Ass}
\DeclareMathOperator{\Min}{Min}
\DeclareMathOperator{\Supp}{Supp}
\DeclareMathOperator{\height}{height}
\DeclareMathOperator{\Hom}{Hom}
\DeclareMathOperator{\ann}{ann}
\DeclareMathOperator{\End}{End}
\DeclareMathOperator{\coker}{coker}
%\DeclareMathOperator{\ker}{ker}
\DeclareMathOperator{\rank}{rank}
\DeclareMathOperator{\im}{im}
\DeclareMathOperator{\Ext}{Ext}
\DeclareMathOperator{\Tor}{Tor}
%\DeclareMathOperator{\dim}{dim}
\DeclareMathOperator{\HH}{H}

\DeclareMathOperator{\lcm}{lcm}

\def\ra{\rightarrow}
\newcommand{\m}{\mathfrak{m}}
\newcommand{\C}{\mathbb{C}}
\newcommand{\Q}{\mathbb{Q}}
\newcommand{\R}{\mathbb{R}}
\newcommand{\N}{\mathbb{N}}
\newcommand{\ov}[1]{\overline{#1}}


\newcommand{\Z}{\mathbb{Z}}
\DeclareMathOperator{\pdim}{pdim}
\DeclareMathOperator{\kos}{kos}
\DeclareMathOperator{\embdim}{embdim}
\DeclareMathOperator{\depth}{depth}
\DeclareMathOperator{\grade}{grade}
\DeclareMathOperator{\injdim}{injdim}

\title{}
\date{\vspace{-0.5in}}

\makeatletter
\g@addto@macro\@floatboxreset\centering
\makeatother

\theoremstyle{definition}
\newtheorem{problem}{Problem}


\begin{document}

\thispagestyle{fancy}
\pagestyle{fancy}
\rhead{UNL}
\lhead{Math 906}
\chead{Final Exam}

\vspace{3em}

\begin{center}
	{\LARGE Final Exam \\}
	Due Friday, May 1, 2026
\end{center}

\

\noindent
{\bf Instructions:}
Write complete and clear solutions to each problem. If you use Macaulay2, email me a~.m2 file with your work, containing comments.
The only allowed resources are me, Macaulay2, and our class notes.
You can work on the problems with your classmates, but only after spending at least $\min \lbrace \text{time it takes you to solve the problem}, \text{1 hour} \rbrace$ per problem. 


\vspace{1em}


\noindent
Turn in {\bf 5 problems} of your choosing. Any problem you do not turn is now a known theorem.

\



\begin{problem}
	Let $(R, \m, k)$ be a Gorenstein local ring. Show that the minimal injective resolution for $R$ has the form
\[ \xymatrix@C=1.3em{
 0 \arrow[r] & \displaystyle\bigoplus_{\height(P) = 0} E(R/P) \arrow[r] & \displaystyle\bigoplus_{\height(P) = 1} E(R/P) \arrow[r] & \cdots \arrow[r] & \displaystyle\bigoplus_{\height(P) = \dim(R)-1} E(R/P) \arrow[r] & E(R/\m) \arrow[r] & 0.
}\]
\end{problem}


\vspace{0.5em}


\begin{problem}
	Let $(R, \m, k)$ be a Cohen-Macaulay ring with a canonical module $\omega_R$. Show that $R$ is Gorenstein if and only if $\omega_R$ is a cyclic $R$-module.
\end{problem}

\vspace{0.5em}

\begin{problem}
Let $(Q, \m, k)$ be a regular local ring, and $R = Q/I$ Cohen-Macaulay with $\pdim_Q(R) = n$.
Show that $R$ is Gorenstein if and only if the betti sequence for $R$ over $Q$ is symmetric, meaning that for all $i$
\[ \beta_i^Q(R) = \beta_{n-i}^Q(R).\]

	\noindent
	Hint: Apply $\Hom_Q(-,Q)$ to the minimal free resolution for $R$ over $Q$.
\end{problem}


\vspace{0.5em}

\begin{problem}
	Let $(Q, \m, k)$ be a regular local ring of dimension $d$ and let $R = Q/I$ with $I \subseteq \m^2$.
		
		\vspace{-0.3em}
	\begin{enumerate}[label=\alph*), itemsep=-0.1em]
		\item Show that if $R$ is a Cohen-Macaulay ring of dimension $d-1$, then $R$ is a hypersurface.
		\item Show that if $R$ is a Gorenstein ring of dimension $d-2$, then $R$ is a complete intersection.
	\end{enumerate}
\end{problem}



\vspace{0.5em}

\begin{problem}
Consider the nonstandard graded ring
\[ R = \mathbb{Q}[t^3, t^{13}, t^{47}].\]

Write an explanation of your calculations in the~.pdf, and include an auxiliary~.m2.
	
	\vspace{-0.3em}
	\begin{enumerate}[label=\alph*), itemsep=-0.1em]
		\item Use Macaulay2 to find a minimal presentation for $R$ of the form $R=\mathbb{Q}[x,y,z]/I$, where $I$ is a homogeneous ideal in $\mathbb{Q}[x,y,z]$ under an appropriate choice of grading.
		\item Can $I$ be generated by a regular sequence?
		\item Is $R$ Gorenstein?
		\item Is $R$ Cohen-Macaulay? Explain your answer with a proof, and confirm it with Macaulay2.
		\item Let $\m = (x,y,z)$. Does $R_\m$ have a canonical module? If so, find a presentation for $\omega_{R_\m}$.
	\end{enumerate}
\end{problem}


\vspace{0.5em}


\begin{problem}
	Let $R$ be a noetherian ring of dimension $1$.
	
	\vspace{-0.3em}
	\begin{enumerate}[label=\alph*), itemsep=-0.1em]
		\item Show that if $R$ is reduced,\footnote{Recall that a ring is reduced if $\sqrt{(0)} = (0)$.} then $R$ is Cohen-Macaulay.
		\item Show that if $R$ is local, then $R$ has a finitely generated MCM module.
	\end{enumerate}
\end{problem}

\vspace{1em}
\end{document}